\section{Related Work}
\label{sec:related_work}

\para{English sycophancy under challenge.} Recent work shows that sycophancy is a stable failure mode in RLHF-era assistants rather than an isolated prompt artifact. \citet{sharma2024sycophancy} document challenge-induced answer reversals across several English tasks and connect the behavior to preference data. \citet{wei2023synthetic} show that instruction tuning and scale can amplify agreement-seeking behavior, while simple synthetic intervention data can reduce it. \citet{chen2024yesmen} add a mechanistic account, arguing that a small subset of components can drive the behavior and that targeted mitigation can suppress it. Our work builds on this line directly, but changes the central axis of variation from prompt type to language resource.

\para{Beyond direct factual agreement.} Sycophancy is broader than agreeing with an explicitly wrong factual claim. \citet{cheng2025elephant} expand measurement to social settings such as validation, framing, and preserving the user's self-image. This broader view is useful because it shows that challenge-turn factual sycophancy is only one part of the phenomenon. We focus on the factual case deliberately: it gives objective labels and lets us ask whether weaker multilingual factual representations make challenge-induced capitulation more likely.

\para{Cross-lingual sycophancy is still underexplored.} Direct multilingual evidence remains sparse. The closest result in our literature review is a recent Hindi Beacon extension that reports larger sycophancy gaps after cultural adaptation than after literal translation \cite{sattigeri2026hindi}. That paper is valuable because it shows that cross-lingual effects can be real, but it is still single-language and uses a different prompt family. Our study instead spans seven languages, uses two benchmarks, and includes both behavioral and intervention analysis.

\para{Multilingual factuality gaps.} A separate literature shows that factual reliability varies substantially across languages. \citet{chataigner2024multilingual} report multilingual hallucination gaps in free-form generation, including weaker performance in several lower-resource settings. \citet{longpre2020mkqa} provide the canonical aligned multilingual factual QA benchmark, which remains useful because it holds answer identity constant across languages. We draw on this literature to test the factual side of the hypothesis, then compare it against false-premise rejection under pressure.

\para{Positioning our work.} Unlike English sycophancy studies, we ask whether resource level changes challenge behavior directly. Unlike multilingual factuality work, we pair language variation with an explicit second-turn challenge. Unlike recent single-language cross-lingual sycophancy studies, we show that the expected inverse-resource trend does not hold cleanly on our benchmark design. The result is less dramatic than the original hypothesis, but it is methodologically important: multilingual challenge evaluations can conflate agreement pressure with conservative refusal.
